\chapter{Conclusion}
Sleep staging is a complex clinical task that has traditionally been performed by 
human experts who meticulously analyze continuous physiological recordings.
Recently, researchers have increasingly utilized deep learning models to automate and scale
this process. In this thesis, we developed a deep neural network based on the Transformer
architecture to classify sleep stages from raw EEG signals. We demonstrated the inherent limitations
of a generalist baseline model that processes data without demographic context, and we
subsequently introduced and evaluated three distinct methodologies for incorporating patient age
to improve generalization and robustness.

Our baseline evaluations revealed that while the unconditioned model performed well on the largest adult dataset,
it struggled to generalize across other developmental stages, exposing two critical vulnerabilities: an over-reliance
on information shared across channels and demographic bias. For underrepresented cohorts, such as children and infants,
the full 6-channel model achieved decent accuracy by exploiting dependencies between channels that did not generalize well.
This became evident when evaluating the model using only 3 channels, which caused predictive accuracy for
these younger cohorts to collapse catastrophically.

Conversely, the model exhibited a different yet equally severe limitation when processing teenage data. When we evaluated the model on all 6 channels, accuracy for teenagers was the lowest among all cohorts; however, accuracy substantially improved when we reduced the spatial dimensionality to 3 channels. 

We argue this anomaly stems from the interaction between extreme dataset imbalance and adolescent neurodevelopment. Because teenagers share fundamental temporal sleep dynamics with adults, the network attempted to map teenage recordings into the dominant adult feature space. However, due to ongoing adolescent brain maturation, teenage EEG exhibits distinct amplitudes and spatial distributions compared to adult sleep. When evaluating across all 6 channels, these spatial differences conflicted with the network's adult-biased spatial filters, acting as noise. 
We hypothesize that reducing the input dimensionality to 3 channels effectively provided architectural regularization: by removing the extra spatial channels, we prevented the network from enforcing rigid adult spatial rules, allowing it to leverage the shared temporal sleep features that generalize across both demographics.

To address these demographic disparities, our first proposed solution 
incorporated the patient's age cohort as an explicit demographic prior 
via an additional embedding layer. By offloading the task of age deduction,
the network was able to learn cohort-specific representations without having
to implicitly extract developmental stages from the raw EEG data. 
While this balancing effect resulted in a marginal decrease in adult accuracy,
performance on the underrepresented and morphologically distinct cohorts improved significantly.
Most crucially, this demographic conditioning served as a powerful 
regularizer against spatial brittleness. When evaluating on only 3 channels,
we successfully mitigated the catastrophic performance collapse we previously observed in the children cohort,
maintaining a significantly higher stability.

Interestingly, we observed an anomaly when we natively trained and evaluated the age-conditioned model
on only 3 channels: accuracy for 
the teenage cohort actually decreased (77.38\%) compared to the 
unconditioned 3-channel baseline (81.28\%). We speculate that this 
drop is a consequence of representational partitioning. In the unconditioned 
3-channel baseline, the absence of age group information allowed teenagers 
to seamlessly leverage the massive volume of shared features learned from the 
adult majority. By explicitly providing an age embedding, the conditioned model 
artificially segregated the teenage representations from the adults. 
Because their core temporal waveforms are highly similar, this 
explicit separation prevented the teenage cohort from fully capitalizing
on the shared adult temporal priors, demonstrating that demographic conditioning 
can occasionally hinder transferability when the underlying domain features are 
already strongly aligned.

To further investigate the impact of demographic variance, our second approach
completely partitioned the dataset to train independent, age-specific models.
Rather than processing all cohorts through shared temporal and spatial weights,
this isolation allowed each network to optimize its feature extraction exclusively
for its target demographic. However, this approach exposed a critical vulnerability
regarding dataset volume. The specialized model for the children cohort, which
possessed the fewest training samples, yielded the poorest overall native
performance (70.37\%). Furthermore, when we subjected this model to the 3-channel missing data
simulation, its accuracy catastrophically plummeted to 34.69\%. This
confirms that in the absence of sufficient training data, the full 6-channel
spatial resolution leads to overparameterization.

Conversely, partitioning the dataset yielded mixed but highly illuminating 
results for the larger cohorts. The isolated adult model actually underperformed 
the generalist baseline (75.74\% vs. 76.76\%), indicating that the baseline network
benefited from the regularization and diversity provided by the aggregate dataset
(which included over 310,000 total samples across all ages). However, the
age-specific approach proved exceptionally beneficial for the teenage cohort,
achieving an impressive 83.38\% accuracy on the full spatial resolution. We attribute
this surge in performance to the network's ability to calibrate its spatial filters
directly to adolescent-specific distributions. Unlike the children
cohort, the teenage dataset was sufficiently large to avoid severe overfitting.
Freed from the dominant adult spatial biases that constrained the baseline and
conditioned models, the isolated teenage network could fully leverage all 6
channels to learn optimized, robust representations unique to adolescent sleep
architecture.

To address the data scarcity we observed in the younger demographics, our third
proposed solution utilized transfer learning. By utilizing the partitioned adult
networks as base models, we aimed to leverage foundational sleep features learned
from the massive adult majority and fine-tune them on the smaller,
underrepresented cohorts. This approach, however, yielded a complex trade-off
between absolute peak performance and diagnostic robustness. When we trained and
evaluated children and infant models on the full 6-channel resolution, they slightly
underperformed their from-scratch, age-specific counterparts. We attribute this
to weight anchoring and negative transfer; we had so deeply optimized the pre-trained spatial filters
for adult neurophysiology that the network struggled to fully
adapt them to the unique spatial topographies of children and infants during fine-tuning.

Despite this slight drop in peak accuracy, transfer learning demonstrated
profound clinical value in mitigating spatial brittleness for data-starved
cohorts. When we subjected the fine-tuned children's model to the simulated loss
of 3 channels, it maintained an accuracy of 57.27\%—a massive 23 percentage point 
improvement over the catastrophic 34.69\% crash we observed in the isolated model. 
By grounding the network in generalized adult priors, the fine-tuning process
successfully prevented the data-scarce children's model from memorizing
brittle spatial shortcuts. Finally, this experiment reinforced our previous findings
regarding architectural regularization: for the infant cohort,
fine-tuning a natively 3-channel model actually outperformed fine-tuning the
6-channel model. By restricting the spatial dimensions during transfer learning,
we successfully stripped away the rigid, adult-specific spatial biases, allowing
the network to cleanly transfer highly generalizable temporal sleep waveforms
across different age stages.

Finally, to contextualize our contributions against established clinical software,
we evaluated our architectures against YASA (Yet Another Spindle Algorithm),
a popular off-the-shelf sleep staging tool. Even under an optimistic evaluation
protocol that granted YASA the optimal demographic assumptions, the model achieved
an average accuracy of 67.60\% for adults, 65.39\% for teenagers, 60.94\% for children,
and a mere 44.96\% for infants. Across every demographic cohort, our proposed
Transformer architectures substantially outperformed this established baseline.
Our unconditioned model achieved higher accuracy on adults (76.76\% vs. 67.60\%),
while our demographic-aware approaches dramatically outpaced YASA on younger
cohorts—most notably for teenagers (reaching 83.38\% vs. 65.39\%) and children
(reaching 76.18\% vs. 60.94\%).

We observed the most profound gap in the infant cohort, where our age-conditioned
and age-specific models outperformed YASA by nearly 30 percentage points (73.79\%
and 74.87\%, respectively, vs. 44.96\%). This dramatic performance collapse in the
established model underscores a critical limitation of traditional off-the-shelf
tools: because they rely on single-channel rules and feature extractors trained
predominantly on adult populations, they fail to generalize to the
pediatric populations. Ultimately, this comparison
validates the core motivation of this thesis, proving that off-the-shelf software is
ill-equipped for non-adult sleep staging and demonstrating the necessity of age-aware
deep learning models.

Through these experiments, this thesis has demonstrated the critical importance
of explicitly accounting for demographic variance in deep learning models
for sleep staging. Furthermore, our findings highlight a significant
methodological risk: relying solely on standard test-set accuracy can yield
misleading conclusions regarding model competence. As we demonstrated, seemingly
minor data changes such as the simulated loss of a few clinical sensors 
can expose severe spatial brittleness and trigger catastrophic 
performance degradation. Consequently, when developing models intended for 
clinical generalization, it is important to subject them to rigorous 
stress-testing across multiple real-world failure scenarios.

While this research focused exclusively on patient age, it is highly probable
that similar representational biases exist across other demographic variables,
such as neurological health conditions or underlying sleep disorders,
which also fundamentally alter EEG morphology. A natural limitation of
this study is the absence of prospective validation in a live clinical
environment. Future work should prioritize deploying these networks alongside
human sleep experts to directly compare the models' generalized representations
against established clinical workflows. Ultimately, however, this thesis
establishes the profound potential of Transformer architectures for
automated sleep staging, proving their high clinical utility when datasets
and training methodologies are carefully designed to balance both accuracy and robustness.